\documentclass[conference]{IEEEtran}
\IEEEoverridecommandlockouts

\usepackage{cite}
\usepackage{amsmath,amssymb,amsfonts}
\usepackage{algorithmic}
\usepackage{graphicx}
\usepackage{textcomp}
\usepackage{xcolor}

\def\BibTeX{{\rm B\kern-.05em{\sc i\kern-.025em b}\kern-.08em
T\kern-.1667em\lower.7ex\hbox{E}\kern-.125emX}}

\begin{document}

\title{Report 1}

\author{
\IEEEauthorblockN{Leander Ziehm}
\IEEEauthorblockA{
Deggendorf Institute of Technology\\
Leander-Arun.ziehm@stud.th-deg.de
}
}

\maketitle



\section{Data Preprocessing}

The project began with an initial data inspection to assess data quality, completeness, and temporal consistency. First, we checked for missing values across all records. No missing values were found, as all cells contained valid entries; there were neither empty strings nor NaN values present in the dataset.

Next, we evaluated duplicate records. A uniqueness check was performed on full rows as well as on subsets of attributes. No duplicate entries were detected when considering all columns, including the \textit{Short Name} identifier. Even after removing the \textit{Short Name} column, no duplicates were observed in the dataset. Only in highly reduced feature configurations (after removing additional distinguishing attributes) could repeated values occur; however, these cases were not considered true duplicates since they represent valid recurring observations in time-series data rather than erroneous redundancy.

To ensure temporal integrity, we further verified the continuity of timestamps. Data was grouped by cell, and for each cell the time series was checked from the minimum timestamp onward. We confirmed that observations were recorded at consistent 15-minute intervals, resulting in a continuous sequence of 5760 timestamps per cell. This ensured that no gaps existed in the time series.

In addition, we verified the uniqueness of $(cell, timestamp)$ pairs to ensure that each measurement was uniquely recorded in time. No duplicate observations were found under this constraint.

\subsection{Normalization / Scaling}

Normalization was evaluated as an important preprocessing step, particularly for machine learning models sensitive to feature magnitude differences. This is especially relevant for distance-based clustering methods such as K-Means, where Euclidean distance can be dominated by features with large numeric ranges.

We considered three scaling methods from the \texttt{scikit-learn} library:

\begin{itemize}
    \item \textbf{MinMaxScaler} – scales features to a fixed range, typically [0, 1], but is sensitive to outliers.
    \item \textbf{StandardScaler} – standardizes features using z-score normalization.
    \item \textbf{RobustScaler} – uses median and interquartile range, making it more robust to outliers and skewed distributions.
\end{itemize}

Given the nature of the dataset, which exhibits bursty and highly skewed network traffic behavior with frequent extreme values, RobustScaler was created to better deal with outliers so it has a might be the best pick but we wanted to also compare it to StandardScaler to see the impact and difference. MinMax scaling was discarded due to its sensitivity to outliers.

To avoid data leakage, scaling parameters were fitted only on the training data and subsequently applied to the test set.

Additionally, we are considering logarithmic transformation as a potential future enhancement to reduce skewness as it helps visually see the distributions so it might also help the models.

\section{Cell Clustering Exploratory Analysis}

Several clustering approaches were explored to identify meaningful structure in the dataset. Due to computational constraints and the exploratory nature of the study, we focused on efficient and interpretable methods.

The following techniques were evaluated:

\begin{itemize}
    \item K-Means Mini-Batch  clustering on raw data
    \item K-Means Mini-Batch with PCA-based dimensionality reduction
    \item UMAP combined with HDBSCAN for density-based clustering
\end{itemize}

For K-Means, we used silhouette score analysis to determine the optimal number of clusters. On raw data, the optimal configuration suggested approximately 2 clusters. However, after feature engineering, the optimal number increased to approximately 8 clusters, indicating richer structure in the transformed feature space.

Based on this, we selected these variables to investigate how the minimum and maximum cluster sizes were affected under different configurations, including feature engineering and the application of different scaling methods. For DBSCAN, the minimum number of points per cluster was initially set to 2 in order to explore the smallest meaningful groupings that could still be identified. Potentially, this threshold could be reduced further to make the AI model more practical by reducing the number of models required per cell cluster. However, at this stage, the main objective was exploratory: to understand how many clusters the method would produce, which ranged between 13 and 17 depending on the configuration. A more detailed analysis of these results is provided in the appendix.

Evaluating the clustering quality proved challenging because all methods were trained in an unsupervised setting. We considered using silhouette score together with additional clustering metrics, although this introduces some limitations since K-Means already optimizes for silhouette score. As a result, we were uncertain whether this alone would provide a fair comparison across methods. To complement the quantitative evaluation, we also developed a dashboard for manual inspection of the clustering outputs, which we plan to expand upon in future work.

Future research could investigate alternative evaluation strategies and more advanced feature representations. In particular, incorporating features that better capture spikes, seasonality, and other complex temporal patterns may improve the clustering quality and better align the results with human intuition and domain-specific expectations.
Although Dynamic Time Warping (DTW) clustering was initially considered due to its suitability for time-series alignment, it was not fully implemented due to time constraints.

Overall, the results indicate that the dataset contains more than two distinct behavioral patterns, and that feature engineering significantly improves cluster separability.

\section{Feature Engineering}

The objective of the feature engineering process was to generate a broad set of temporal and statistical features and later allow feature selection methods such as forward selection, backward elimination, correlation analysis, or model-based importance techniques to identify the most relevant features.

We considered which features to include. The first idea was to add information about the day of the week using numeric encoding. A column was added where Monday = 0 and Sunday = 6, allowing the model to better interpret weekly patterns. To improve temporal representation, the date was also split into separate components.

In order for the model to better capture periodic patterns, time-based features such as hour, day of week, and month were introduced. For temporal variables, cyclical encoding using sine and cosine transformations was applied to preserve their periodic nature. 

To capture changes over time, difference features such as \texttt{diff\_1} and \texttt{diff\_24} were created, where \texttt{diff\_1} refers to the difference at the previous timestep (i.e., 15 minutes). Lag features were also introduced to model temporal dependencies and historical behavior.

To smooth variability and capture different time-scale behaviors, rolling statistics such as rolling mean and rolling standard deviation were used over multiple windows.

Since spikes often occur before holidays, weekends, or during weekends, calendar-based indicators such as \texttt{is\_holiday}, \texttt{is\_preholiday}, and \texttt{is\_weekend} were included to account for behavioral changes.

Correlation analysis using both Pearson and Spearman methods was performed to evaluate linear and non-linear relationships in the dataset, both before and after feature engineering. This was important since network traffic behavior may increase non-linearly with user activity during congestion periods. Overall, the engineered features were designed to support both clustering and time-series analysis tasks.

\subsection{Final Feature Set}

The final feature set implemented in \texttt{feature\_engineering.py} includes:
\begin{itemize}
    \item Temporal features: \texttt{hour}, \texttt{minute}, \texttt{dayofweek}, \texttt{day}, \texttt{month}, \texttt{weekofyear}, \texttt{timeslot}
    \item Cyclical encodings: \texttt{hour\_sin}, \texttt{hour\_cos}, \texttt{dow\_sin}, \texttt{dow\_cos}, \texttt{timeslot\_sin}, \texttt{timeslot\_cos}
    \item Difference features: \texttt{diff\_1}, \texttt{diff\_24}
    \item Lag features: \texttt{lag\_1}, \texttt{lag\_4}, \texttt{lag\_12}, \texttt{lag\_24}, \texttt{lag\_48}, \texttt{lag\_96}, \texttt{lag\_192}, \texttt{lag\_288}, \texttt{lag\_672}
    \item Rolling statistics: \texttt{roll\_mean\_4}, \texttt{roll\_std\_4}, \texttt{roll\_mean\_12}, \texttt{roll\_std\_12}, \texttt{roll\_mean\_24}, \texttt{roll\_std\_24}, \texttt{roll\_mean\_96}, \texttt{roll\_std\_96}
    \item Calendar indicators: \texttt{is\_holiday}, \texttt{is\_preholiday}, \texttt{is\_weekend}
\end{itemize}

\section{Future Work}

Future work will extend temporal feature engineering by incorporating long-range lag dependencies, starting with a 1-week lag feature to capture recurring weekly traffic behavior. Further investigation will be conducted to determine whether 2-week and 3-week lag structures provide additional predictive value or introduce redundancy.

Additionally, advanced statistical and behavioral descriptors will be explored to better characterize network traffic dynamics, including:

\begin{itemize}
    \item Kurtosis for identifying heavy-tailed or extreme traffic fluctuations
    \item Daily similarity scores to quantify recurring daily usage patterns
    \item Burstiness metrics to capture sudden traffic surges and irregular activity
    \item Shannon entropy to measure traffic randomness and variability
    \item Gini coefficient to evaluate distribution imbalance and traffic concentration over time
\end{itemize}

These higher-order temporal and distribution-based features will be evaluated based on their contribution to clustering quality, anomaly detection sensitivity, and forecasting performance in telecom traffic analysis.
\clearpage
\section{Appendix}

Cell Clustering Exploratory Analysis Table:
\begin{table}[ht]
\centering
\begin{tabular}{l c c c c c c}
\hline
Method & Enabled Feature Engineering & Used Standard Scaler & n\_clusters & Min Cluster Size & Max Cluster Size & Mean Cluster Size \\
\hline
MiniBatchKMeans\_raw & True & True & 2 & 57.0 & 169.0 & 113.0 \\
PCA\_KMeans & True & True & 2 & 59.0 & 167.0 & 113.0 \\
MiniBatchKMeans\_raw & True & False & 2 & 1.0 & 225.0 & 113.0 \\
PCA\_KMeans & True & False & 2 & 32.0 & 194.0 & 113.0 \\
MiniBatchKMeans\_raw & False & True & 2 & 58.0 & 168.0 & 113.0 \\
PCA\_KMeans & False & True & 2 & 60.0 & 166.0 & 113.0 \\
MiniBatchKMeans\_raw & False & False & 2 & 1.0 & 225.0 & 113.0 \\
PCA\_KMeans & False & False & 2 & 44.0 & 182.0 & 113.0 \\
MiniBatchKMeans\_raw & True & True & 8 & 1.0 & 50.0 & 28.2 \\
PCA\_KMeans & True & True & 8 & 3.0 & 126.0 & 28.2 \\
MiniBatchKMeans\_raw & True & False & 8 & 1.0 & 111.0 & 28.2 \\
PCA\_KMeans & True & False & 8 & 1.0 & 114.0 & 28.2 \\
MiniBatchKMeans\_raw & False & True & 8 & 1.0 & 115.0 & 28.2 \\
PCA\_KMeans & False & True & 8 & 2.0 & 113.0 & 28.2 \\
MiniBatchKMeans\_raw & False & False & 8 & 1.0 & 104.0 & 28.2 \\
PCA\_KMeans & False & False & 8 & 1.0 & 118.0 & 28.2 \\
UMAP\_HDBSCAN & True  & True  & 13.0 & 3.0 & 92.0 & 17.4 \\
UMAP\_HDBSCAN & True  & False & 14.0 & 2.0 & 74.0 & 16.1 \\
UMAP\_HDBSCAN & False & True  & 13.0 & 2.0 & 90.0 & 17.4 \\
UMAP\_HDBSCAN & False & False & 17.0 & 2.0 & 65.0 & 13.3 \\
\hline
\end{tabular}
\end{table}




\end{document}